% Figaro Paper I: The Ungoverned Substrate
% Network States, Blockchain Governance,
% and the Option Both Programs Presuppose Away
% (Companion to Paper H; engages Srinivasan and the De Filippi /
% COALA blockchain-governance research program.)
% Preprint, April 2026

\documentclass[11pt,a4paper]{article}

\usepackage[utf8]{inputenc}
\usepackage[T1]{fontenc}
\usepackage{amsmath,amssymb,amsthm}
\usepackage{mathtools}
\usepackage{booktabs}
\usepackage{graphicx}
\usepackage{hyperref}
\usepackage{cleveref}
\usepackage{enumitem}
\usepackage{xcolor}
\usepackage[margin=1in]{geometry}
\usepackage{natbib}
\bibliographystyle{plainnat}

\theoremstyle{remark}
\newtheorem{remark}{Remark}[section]

\newcommand{\figaro}{\textsc{Figaro}}
\newcommand{\fig}{\textsc{Fig}}

\title{%
  The Ungoverned Substrate:\\
  Network States, Blockchain Governance,\\
  and the Option Both Programs Presuppose Away%
}

\author{
  Alessandro Daliana\thanks{Corresponding author. \texttt{adaliana@figaro.org}}
}

\date{April 2026}

\begin{document}

\maketitle

% ════════════════════════════════════════════════════════════════════
\begin{abstract}
Two contemporary research programs have addressed the question of
what kind of polity a network-native coordination architecture
might support. The first, articulated most prominently by
\citet{srinivasan2022network}, proposes the \emph{network state}:
a community that forms online, accumulates social cohesion and
shared institutions, and ultimately seeks territorial and
diplomatic recognition as a sovereign-equivalent entity. The
second, developed by \citet{defilippi_wright_2018} and the
collaborators around the Coalition of Automated Legal
Applications (COALA), examines the legal and governance
consequences of running rules through code: \emph{lex
cryptographica}, decentralized autonomous organizations,
on-chain governance, and the legal scaffolding such arrangements
require to operate alongside conventional jurisdictions. The two
programs differ in goal --- one constitutional, the other
infrastructural --- but share a structural premise.

The shared premise is that the alternative-to-state-coordination
must itself be a \emph{governed entity}. The network state
governs through a constitutional process inherited from the
modern sovereign-state form. The blockchain-governance program
governs through on-chain voting, off-chain coordination, and
hybrid legal-computational arrangements. Both programs treat
governance as a question that arises at the network or protocol
layer and must be answered there. Both programs accordingly
spend substantial effort on the design and analysis of governance
mechanisms appropriate to their respective forms.

Behind both programs sits a longer trajectory in which the rules
of social coordination have been sustained by two mechanisms:
coercion through the sovereign's enforcement apparatus, and
legitimacy by belief sustained by trusted institutions.
Governance, in both senses, is the maintenance work those
mechanisms require to remain in force.

This paper observes that the bonded settlement primitive
examined in the corpus's companion mechanism, institutional, and
political-economy papers admits a third option that neither
program names: an \emph{ungoverned substrate} on which arbitrary
governed entities can be constituted, but which itself is not a
governed entity at all. The substrate does not vote; it does not
upgrade; it does not have a constitution; it has no
representatives, no executive, no judiciary, no protocol-level
discretion of any kind. It executes a small set of invariant
rules and admits arbitrary higher-layer arrangements to be
composed on top of it. Network states can be constituted on the
substrate; so can decentralized autonomous organizations; so
can cooperatives, peer networks, individual sole traders,
ad-hoc agent coalitions, and patterns the substrate's authors
cannot anticipate. The substrate's freedom from governance is
what enables the diversity of governance forms above it.

The substantive claim of this paper is therefore not that
network-state or blockchain-governance programs are wrong on
their own terms ---- both are coherent research programs
addressing real questions ---- but that the architectural option
they presuppose away is the option the bonded primitive realizes.
Recognizing this option does not adjudicate among governed
arrangements; it relocates the governance question from the
substrate, where both programs place it, to the graph, where the
substrate admits it to be answered in arbitrarily many ways. That
relocation is the political-architectural content of the
substrate change with respect to network polity, and it is the
subject of this paper.
\end{abstract}

\medskip
\noindent\textbf{Keywords:}
network state, blockchain governance, lex cryptographica,
ungoverned substrate, decentralized autonomous organization,
infrastructural literacy, governance at the graph

% ════════════════════════════════════════════════════════════════════
\section{Introduction}
\label{sec:intro}

This paper sits between political philosophy, the contemporary
literature on blockchain governance, and the analysis of a
particular settlement primitive. Its starting point is a question
that the companion political-philosophy paper on coercion left
open: when the state's sovereign-coercive function is partially
substituted by bonded pre-commitment in bilateral commerce, what
follows for the proposals that have emerged to organize
coordination at the network layer? The question is sharp enough
to deserve dedicated treatment.

Two such proposals have organized the recent discussion. The
first is the network-state thesis, advanced most prominently by
\citet{srinivasan2022network} and adjacent commentators. The
second is the blockchain-governance research program, developed
by \citet{defilippi_wright_2018} and collaborators around COALA
and parallel academic and practitioner communities. The two are
distinct: the first is normative-prescriptive (a program for
constituting alternative polities), the second is
descriptive-analytical (an account of what blockchain
arrangements imply for law and governance). Both, however, share
a structural assumption that the present paper identifies and
argues is not forced.

The shared assumption is that the network-layer alternative to
state coordination must itself be a \emph{governed entity}. The
network state requires a constitutional founder, a community
self-conception, an explicit governance process, and ultimately
sovereign-equivalent diplomatic standing
(\citealt{srinivasan2022network}). The blockchain-governance
program treats the governance of on-chain arrangements as the
central problem of the field: how decisions are made about the
protocol's evolution, who has voting power, what counts as
legitimate amendment, how governance failures are detected and
corrected (\citealt{defilippi_mannan_reijers_2022}). In both
programs, the substrate of network-native coordination is
itself something that must be governed, and the design and
analysis of that governance is the central work.

The bonded settlement primitive examined in the corpus's
companion mechanism paper provides a different option that
neither program names: an \emph{ungoverned substrate}. The
kernel is invariant. It admits no upgrades, no proposals, no
votes, no protocol-level discretion. There is no admin, no
owner, no fee parameter to vote on, no governance token to
hold. There is no entity of which the kernel is the
constitution, because the kernel is not a constitution of
anything. It is a small set of rules executed on a public
infrastructure, on top of which arbitrary governed entities
(network states, decentralized autonomous organizations,
cooperatives, peer networks, sole traders, agent coalitions)
can be composed --- or which bilateral parties can use without
constituting any entity at all.

\paragraph{The thesis.}
The thesis is structural. (1) Both the network-state program and
the blockchain-governance program assume that the alternative to
state coordination must be a governed network entity. (2) The
assumption is not architecturally forced: a coordination
substrate can be invariant rather than governed, and the
decision to make it one or the other is a design choice. (3) The
ungoverned-substrate option, where it is taken, does not abolish
governance --- it relocates governance from the substrate layer
to the graph layer, where any pattern of governance (or absence
of governance) can be composed. (4) The relocation is the
political-architectural content of the substrate change, and it
clarifies what each program is actually proposing: the network
state is one possible graph, blockchain-governance arrangements
are another, and neither is the substrate.

\paragraph{What this paper is and is not.}
This paper is a short essay in the political architecture of
network coordination. It does not extend the underlying
mechanism, the institutional argument, or the political-economy
diagnostic of the companion papers, nor does it extend the
political-philosophy of coercion treated in the companion paper
on the coercion variable. It engages two contemporary research
programs to clarify what the substrate change implies and does
not imply for the proposals those programs have advanced.

The paper takes no position on whether network states are
desirable, whether DAOs are well-governed, whether on-chain
voting is preferable to off-chain coordination, whether
\emph{lex cryptographica} should be recognized by conventional
jurisdictions, whether COALA's model legal frameworks are
correctly drafted, or whether any other normative or
prescriptive question of network polity has the right answer.
Those are real questions, and the paper's silence on them is
intentional: the substrate change is a precondition for those
questions to be posed coherently at the graph layer, not an
adjudication of how they should be answered.

\paragraph{Position relative to the corpus.}
This paper extends the political-philosophy treatment of the
companion paper on coercion. That paper identified coercion as
a substrate variable: structurally load-bearing in the dominant
tradition's account of bilateral commercial enforcement,
substituted by pre-commitment when the bonded primitive is the
operable architecture. The present paper identifies governance
as a related substrate variable: structurally load-bearing in
both contemporary network-polity programs, substituted by
\emph{nothing} (or, equivalently, by an invariant kernel) when
the bonded primitive is the operable architecture, and relocated
to the graph for arrangements that require it. The two papers
together complete the substrate-relocation argument for the
political-philosophy cluster, as the institutional and
political-economy companions did for the firm and the
left/right partition.

% ════════════════════════════════════════════════════════════════════
\section{The Network-State Program and Its Coordination Gap}
\label{sec:network-state}

We begin with the network-state thesis, both because it is the
more public of the two programs and because the companion
political-economy paper on hegemony identified its specific
unsolved sub-problem as one the bonded primitive resolves.

\paragraph{The thesis as articulated.}
\citet{srinivasan2022network} proposes the network state as a
sequence: an online community forms around a shared interest,
acquires social cohesion through ongoing interaction, develops
shared institutions (norms, currencies, dispute mechanisms),
attains a population large enough to matter and a treasury
large enough to act, acquires territorial presence (initially
distributed, eventually contiguous), and ultimately seeks
diplomatic recognition as a sovereign-equivalent entity. A
distinctive structural feature of the proposal is the
\emph{one commandment}: the network state is to be organized
around a single moral innovation that distinguishes it from the
existing political environment, and on which the founding
community has already converged. The one commandment is the
substrate of cohesion --- the reason members are aligned with
each other before any of the institutional apparatus is built ---
and the network state's institutions then articulate that
commandment into a practice. The sequence is presented as a
contemporary updating of Hirschman's framework on \emph{exit} as
a response to institutional dissatisfaction
(\citealt{hirschman1970exit}): when voice within incumbent states
becomes ineffective, exit to a new entity becomes the alternative,
but exit requires a destination, and the network state is the
destination's design.

The thesis takes governance as constitutive. A network state is
not a network state without explicit shared institutions,
without a self-conception articulated by a founder or coalition,
without a process for collective decision, without internal
adjudication of disputes, and ultimately without external
recognition by other sovereigns. Each of these is a governance
function, and the design of each is part of what the
network-state program addresses.

\paragraph{The unsolved sub-problem.}
The thesis names an ``economic engine'' as a constitutive stage
between community formation and recognition --- the network state
must be capable of supporting economic activity within its
participants without routing the activity through incumbent
state-clearable infrastructure --- but treats that stage as
relatively under-specified. In practice, commerce within
network-state communities continues to route through fiat rails,
incumbent platforms, and the legal apparatuses of the
participants' physical jurisdictions. The companion
political-economy paper on hegemony noted that this is not a
minor implementation detail: the missing economic engine is the
substrate gap, and the bonded primitive supplies it.

The supply is structural. Bilateral commercial coordination
between network-state participants, between participants and
non-participants, and between two non-participants who happen to
share the substrate, can be conducted through bonded commitments
that depend on no incumbent state apparatus and no
network-state-internal governance. The primitive does not
constitute the network state; it is the substrate on which
network states (and other arrangements) can be constituted
without retaining their dependency on the very state apparatuses
they are exiting from.

\paragraph{What the thesis does not require.}
A subtler observation is that the substrate the network-state
program needs does not have to be one the network state
\emph{governs}. Treating the substrate as something the network
state should constitute under its own governance --- a
network-state-native blockchain, a network-state-issued currency
governed by the network state's process, a network-state-operated
arbitration system --- recreates within the network state the
very substrate-governance problem the network state was supposed
to resolve relative to the incumbent state. The substrate the
network state actually requires is one that no entity governs,
which any entity (the network state, the incumbent state, a
private cooperative, a single trader) can use, and whose
neutrality is what makes it suitable for cross-affiliation
coordination.

The one-commandment structure makes this point sharper. The
network state's coordination problem is not how to make
participants who already share a moral commitment cooperate ---
that is what the one commandment provides --- but how to give
the cooperation an enforceable form that does not import either
the incumbent state's apparatus or a network-state-internal
sovereign. The bonded primitive supplies that enforceable form
without supplying any moral content: it does not say what the
participants should agree about, and it does not produce a
commandment of its own. A network state organized around any
one commandment can use the substrate without the substrate
constraining what the commandment can be. The neutrality is
not a deficit relative to the network-state program; it is the
property the program needs to avoid recapitulating
substrate-governance at the network-state level.

This is the option the network-state program presupposes away
when it treats every functional layer as something the network
state must eventually govern. The bonded primitive's substrate
is not a layer of the network state; it is a layer beneath the
network state and beneath the incumbent state alike, on which
both can compose, and which is not the property of either.

\begin{remark}[On exit and the missing destination]
Hirschman's analysis of exit assumes that exit is meaningful
only when it has a destination. The network-state program is
constructive about destinations: it specifies what the
destination should look like (a sovereign-equivalent network
entity) and what work is required to constitute it. The bonded
primitive contributes to the analysis of exit a different
observation: exit can be partial. A participant can exit the
incumbent state's commercial-coordination apparatus (by using
the bonded substrate for commerce) without exiting the
incumbent state's other functions (citizenship, criminal-law
protection, public goods, defense). Partial exit was difficult
under prior architectures because the incumbent state's
commercial-coordination apparatus was bundled with its other
functions. The substrate change unbundles. Whether the
participant goes on to constitute a network state is a separate
choice that the substrate does not make for them.
\end{remark}

% ════════════════════════════════════════════════════════════════════
\section{Blockchain Governance and the Confidence Machine}
\label{sec:de-filippi}

We turn now to the blockchain-governance research program, which
addresses the same substrate question from a different angle and
arrives at a more textured account of what governing
network-native coordination involves.

\paragraph{The program in outline.}
\citet{defilippi_wright_2018} examine the legal consequences of
running rules through code: \emph{lex cryptographica} as an
emerging body of computational law, smart contracts as
self-executing agreements, decentralized autonomous organizations
as code-mediated coordination structures, and the
hybrid arrangements that emerge when computational rules
encounter conventional jurisdictions. The framing extends the
``code is law'' formulation of \citet{lessig2006code} into a
positive program: where Lessig observed that code increasingly
performs the regulatory function law had historically performed,
the De~Filippi--Wright account asks what it would mean for code
to become recognized law in its own right, with its own forms,
its own jurisprudence, and its own integration with conventional
legal systems. The work has been
extended by parallel research on
on-chain governance design, the legal recognition of DAOs in
specific jurisdictions, model legal frameworks (developed by
COALA and adopted in modified form by several U.S.~state
legislatures), and the relationship between code-based
arrangements and conventional legal institutions.

A consistent thread across this work is that the apparent
elimination of governance through code is illusory. Rules are
written by people; rules require interpretation when situations
arise that the rules' authors did not anticipate; rule changes
are required when external conditions shift; and the
infrastructure on which the rules run (validators, sequencers,
exchanges, fiat on-ramps) introduces governance through the back
door even when the on-chain protocol claims to have eliminated
it. The blockchain-governance program's central observation is
that code-mediated coordination does not abolish governance --- it
relocates and transforms governance, often in ways that are less
visible and accountable than the conventional institutions
governance has historically inhabited.

\citet{walch2017path, walch2019deconstructing} sharpens this
critique. \citeauthor{walch2017path} argues that the language
through which the blockchain field describes itself --- in
particular the words ``decentralized,'' ``trustless,'' and ``code
is law'' --- functions as a discourse that obscures rather than
describes the arrangements at issue: ``decentralized'' protocols
typically depend on small core-developer groups whose decisions
carry enormous practical weight, and the relationship between
those developers and the protocol's participants is properly
characterized as \emph{fiduciary} even though no formal fiduciary
duty has been recognized.
\citeauthor{walch2019deconstructing} extends the argument by
deconstructing ``decentralization'' as a unitary property:
distinct dimensions (architectural, political, logical, economic)
can vary independently, and a protocol can be decentralized along
some dimensions while being highly centralized along others. The
hidden-fiduciary critique is the strongest available challenge to
any claim that a protocol has eliminated governance.

\citet{werbach2018blockchain} arrives at an adjacent point from
the trust-theory side: blockchain protocols substitute one
\emph{architecture of trust} for another, and the new architecture
has its own institutional preconditions (mathematical proof,
adversarial review, infrastructural reliability) that are no less
demanding than the trust apparatuses they displace --- only
differently distributed.

\paragraph{The confidence machine.}
\citet{defilippi_mannan_reijers_2022}, building on adjacent work
by \citet{reijers_coeckelbergh_2018} that read blockchain as a
narrative-technology shaping its participants' conceptions of
trust and value, sharpen this observation
into the framing of blockchain as a \emph{confidence machine}:
an architecture that shifts the locus of trust from people and
institutions to code and mathematics, but that does not
eliminate trust or governance. Confidence in code is required;
the code is written and maintained by people; the people are
governed by their own arrangements (corporate, foundational,
informal); the infrastructure on which the code runs introduces
further dependencies; and the participants' \emph{infrastructural
literacy} ---- their capacity to evaluate the arrangements they
are participating in --- is generally inadequate to the
arrangements' complexity.

The confidence-machine framing is important and largely correct
for the class of arrangements it addresses, which is the class
of \emph{governance-rich} blockchain protocols: systems whose
protocol-level rules are intended to evolve over time, whose
parameters are subject to vote, whose treasuries fund ongoing
development, and whose legitimacy depends on participants'
continued endorsement of the governance process. For such
systems, the confidence-machine analysis exposes the
governance-of-the-governance problem clearly: trust shifted from
human institutions to code does not arrive at a final resting
place; it shifts to the human institutions that govern the code.

\paragraph{The COALA legal-scaffolding program.}
A complementary thread of work, also developed substantially
within COALA, has produced model legal frameworks for DAOs,
templates for the integration of on-chain arrangements with
conventional legal personality, and analyses of how
code-mediated arrangements can be made compatible with the
expectations of courts, regulators, and counterparties operating
under conventional legal systems
(\citealt{defilippi_hassan_2016}). The work is collaborative
rather than antagonistic: it treats code-mediated arrangements
and conventional legal arrangements as complementary rather
than alternatives, and seeks to provide the institutional
scaffolding under which the two can coexist.

This is rigorous and useful work. It identifies that DAOs
constituted as code-mediated entities still need legal
recognition to enter into contracts with non-DAO counterparties,
to hold property in conventional jurisdictions, to limit
liability for participants, and to engage with regulators on
terms regulators can recognize. The COALA model laws are
attempts to make this scaffolding available in standardized
form, so that DAOs need not each negotiate their legal status
ad hoc.

\begin{remark}[On the scope of the confidence-machine analysis]
The confidence-machine analysis applies to the class of
arrangements for which it was developed: governance-rich
blockchain protocols whose parameters and rules evolve under
explicit governance processes. It does not, on its own terms,
require that all blockchain arrangements be of this kind. The
analysis is silent on what happens when an architecture is
\emph{not} governance-rich --- when its rules are invariant,
when no parameter is subject to vote, and when no treasury
funds the protocol's evolution. The present paper does not
contradict the confidence-machine analysis; it identifies a
class of arrangements outside its analytical scope.
\end{remark}

\paragraph{Protocol-layer governance versus substrate-layer governance.}
The confidence-machine framing raises a question that the
ungoverned-kernel claim must answer head-on: even an invariant
kernel runs on infrastructure that is not invariant. Ethereum
validators, sequencers and rollup operators, RPC providers,
exchanges that price the bond-denomination currency, and the
fiat on-ramps through which capital reaches the bonded substrate
are all governed entities, and decisions taken in those layers
affect what happens at the kernel.
\citet{walch2017path}'s hidden-fiduciary critique is the
strongest version of this concern: even if no fiduciary
relationship is recognized at the protocol layer, the
core developers, validator set, and infrastructure providers are
in fiduciary-shaped positions relative to participants who
depend on them.

We accept the concern and locate the kernel relative to it
explicitly. The ungoverned-substrate claim of this paper is a
claim about the \emph{kernel} as a unit of governance, not about
the entire stack on which a participant's transaction depends.
The kernel composes two mechanisms (asymmetric bonding and buyer
dominance with atomic resolution) and admits no decision-position
from which kernel-level governance could be exercised.
\emph{Substrate-level} governance --- of the chain on which the
kernel runs, of the validator set, of the infrastructure
operators --- is not eliminated; it is delegated to the chain's
own governance arrangements (which are, in Ethereum's case,
explicitly informal and exposed for what they are). The
participant's confidence in the kernel decomposes into confidence
in the kernel's bytecode (formally verified, immutable),
confidence in the chain's continued operation
(governance-dependent in the way \citeauthor{defilippi_mannan_reijers_2022}
describe), and confidence in the surrounding infrastructure
(governance-dependent and varying). Only the first is what the
present paper claims is ungoverned. The other two are real
dependencies that the participant inherits from the choice of
underlying chain, and we recommend they be evaluated as such
rather than absorbed into the kernel claim.

% ════════════════════════════════════════════════════════════════════
\section{The Common Premise: Governance at the Substrate Layer}
\label{sec:common-premise}

We now identify the structural premise that the network-state
program and the blockchain-governance program share, and that
the present paper argues is not architecturally forced.

\paragraph{What both programs assume.}
The network-state program assumes that the alternative to state
coordination is a network polity that must be governed: it has
a founder, a constitution, a process for collective decision, a
treasury managed by some authority, internal adjudication, and
external diplomatic standing.

The blockchain-governance program assumes that the alternative
to platform-mediated coordination is an on-chain protocol whose
rules and parameters must be governed: it has a governance
token, a voting process, a developer foundation, a process for
upgrades, and external legal scaffolding to integrate with
conventional jurisdictions.

In both cases, the alternative arrangement is itself a governed
entity. The shape of the governance differs --- constitutional
in the network-state program, computational-democratic in the
blockchain-governance program --- but the structural location of
the governance function is the same: at the network or protocol
layer, where decisions about the system's rules, parameters,
and evolution are made. Both programs accordingly devote
substantial effort to the design and analysis of governance
mechanisms appropriate to their respective forms.

\paragraph{Why the assumption seems necessary.}
The assumption is not arbitrary. Most blockchain protocols of
which the contemporary research community is aware are
governance-rich: their parameters change over time, their rules
admit upgrade paths, their treasuries fund ongoing development.
The largest and most-studied protocols (Ethereum, Bitcoin
through its informal but real social-coordination process,
Tezos, Cosmos chains, Optimism, Maker) all have governance
processes of varying degrees of formality. The
blockchain-governance program emerged as a rigorous account of
these processes, and the network-state program emerged in
parallel to address the related question of what such processes
imply for the political organization of the communities they
govern. The assumption that network-native coordination requires
governance reflects the architectures the programs were
analyzing.

\paragraph{Why the assumption is not architecturally forced.}
A protocol's parameters and rules need not evolve. A coordination
architecture can be designed such that its small invariant kernel
is fixed at deployment, with all variability pushed to layers
above the kernel that the kernel does not constrain. The bonded
settlement primitive examined in the corpus's companion mechanism
paper is such an architecture. The kernel has no admin, no owner,
no upgrade path, no governance token, no treasury, no parameter
that can be voted on, no discretion of any kind. It executes a
small fixed set of rules. Whatever assemblies are constituted on
top of it inherit no governance dependency on the kernel and no
obligation to participate in the kernel's evolution because the
kernel does not evolve.

This is not an attempt to evade governance; it is a
demonstration that governance is a substrate-design choice, not
a substrate-design necessity. The choice has consequences. A
governance-rich protocol can adapt to changing conditions; an
invariant kernel cannot. A governance-rich protocol can correct
errors in its rules; an invariant kernel cannot. A
governance-rich protocol provides governance scaffolding by
default; an invariant kernel provides none. These are real
costs and benefits and the choice between them is substantive.
But the choice exists, and pretending it does not is the move
both programs make when they treat governance at the substrate
layer as constitutive of the kind of arrangement they are
analyzing.

\begin{remark}[On the relationship between governance and adaptability]
The most common defense of substrate-layer governance is that
adaptability is required: external conditions change, errors
need correction, opportunities require parameter updates. This
defense is forceful for protocols whose function is itself
adaptive (lending markets that must respond to market
conditions, scaling layers that must respond to throughput
demand, identity systems that must respond to evolving
identity standards). It is weaker for protocols whose function
is structurally invariant. A protocol that does one well-defined
thing (enforce a bonded bilateral commitment) can be invariant
in a way that a protocol that does many things subject to
external state cannot. The adaptability defense is therefore
specific to function, not generic to protocols.
\end{remark}

% ════════════════════════════════════════════════════════════════════
\section{The Ungoverned Substrate: A Structural Alternative}
\label{sec:ungoverned-substrate}

We now describe the option that the bonded primitive realizes
and that neither contemporary program names.

\paragraph{What ``ungoverned'' means here.}
The kernel is governed in no respect that the contemporary
governance literature would recognize. There is no proposal
process, no voting mechanism, no quorum, no veto, no executive,
no foundation that controls protocol evolution, no developer
group with privileged update access, no upgrade path of any
kind. The deployed bytecode is invariant: identical input
produces identical output for as long as the underlying chain
operates. To change the kernel's behavior requires deploying a
\emph{different kernel} at a different address, which constitutes
a different protocol that the existing one has no relationship
to.

This is not an absence of governance through inattention. It is
governance designed out of the substrate. The kernel's invariants
are fixed by mathematics, formal verification, and code, not by
the consent of an evolving body of participants. The kernel
does not require governance because it does not have decisions
to make. The non-discretion is architectural: the kernel composes
two mechanisms --- asymmetric bonding (Mechanism~1) and buyer
dominance with atomic resolution (Mechanism~2) --- and Mechanism~2
is what eliminates the resolution-time discretionary surface that
governance would otherwise be needed to occupy. Only the buyer can
trigger resolution; resolution settles every active order in the
process simultaneously or not at all; no third party --- arbitrator,
oracle, jury, foundation --- is in a position to defer, partial-resolve,
selectively enforce, or override. The ungoverned-substrate claim
is therefore not a claim that the kernel \emph{happens} to be
small enough to verify; it is a claim that Mechanism~2 leaves
no decision-position from which governance could be exercised
even if a governance body were stood up adjacent to the kernel.

\paragraph{What this requires of the kernel.}
The ungoverned-substrate option is not free. It requires that
the kernel's function be specifiable in advance, that the
specification be small enough to be formally verified, and that
the resulting bytecode be small enough to be auditable to a
standard that warrants trust without ongoing governance. The
companion implementation paper documents how this is achieved
for the bonded primitive: a kernel of approximately two hundred
lines of Solidity, formally verified by multiple independent
verification stacks (TLA$^+$, symbolic execution, property-based
fuzzing, deductive verification), with a small public surface
and no admin functions. The choice to have an ungoverned
substrate is purchased by accepting the constraint that the
substrate's function be amenable to this discipline.

\paragraph{This is not ``code is law.''}
The Lessig formulation \citep{lessig2006code} that ``code is
law'' has been useful as a slogan and damaging as a method.
Useful, because it correctly observed that software increasingly
performs the regulatory function law had historically performed.
Damaging, because the slogan obscures the substantial differences
between what software and law do --- a point both
\citet{walch2017path} and \citet{defilippi_wright_2018} have
emphasized in different registers. Law is interpretively
flexible, jurisdictionally bounded, amendable through legislative
or judicial process, and applied by humans; code is interpretively
rigid, globally identical, immutable except through redeployment,
and applied by execution. Treating one as the other in either
direction produces category errors.

The ungoverned-substrate claim of this paper does not depend on
``code is law.'' It depends on a narrower claim: that the
\emph{settlement primitive} --- a single bilateral function with a
deterministic outcome --- is exactly the class of arrangement for
which formal specification, verification, and immutable
deployment are appropriate, and that this class is much smaller
than ``everything regulated by software.'' Law remains where
law has always been: in the off-chain forums that adjudicate
disputes under conventional doctrines (the companion
evidence-and-law paper develops this in detail), in the
jurisdictional environments within which the parties operate,
and in the interpretive work that any contract --- whether
recorded on chain or off --- ultimately requires. The kernel is
not law; it is settlement infrastructure that produces evidence
on which law operates. The ``code is law'' frame would conflate
the two; we are at pains to keep them distinct.

Most blockchain protocols that have been built do not satisfy
the constraint. Lending protocols, scaling solutions, identity
systems, prediction markets, and many other categories of
on-chain arrangement involve interactions with external state
that the protocol must respond to and that no fixed code can
anticipate. For those, governance is genuinely required because
adaptability is genuinely required. The bonded settlement
primitive is in a different category: the primitive's function
(enforce a bonded bilateral commitment between consenting
parties) is structurally amenable to fixed specification, and
that is what makes the ungoverned-substrate option available
for it.

\paragraph{The composition admit.}
A consequence of the ungoverned-substrate design is that the
substrate admits arbitrary higher-layer arrangements without
constraining them. A community can constitute a network state on
top of the substrate, with whatever governance the network state
chooses (constitutional, computational-democratic, charismatic,
inherited, none). A group of participants can constitute a DAO,
with whatever on-chain or off-chain governance the DAO chooses.
A cooperative can compose itself, with whatever cooperative
governance is appropriate. A peer network can operate without
constituting any entity at all. Two strangers can transact once
and never again. The substrate accommodates each of these
because it makes no commitments about what the parties operating
on it should be.

This is the central architectural observation. The substrate's
ungoverned character is what makes it neutral with respect to
the governance forms that compose on top of it. A
governance-rich substrate would import its governance into
every arrangement constituted on top of it, because the
arrangements would inherit the substrate's governance
dependencies. An ungoverned substrate imports nothing; the
arrangements that compose on top of it owe the substrate
nothing except the use of its public infrastructure, and they
are free to govern (or not govern) themselves in whatever way
their participants choose.

\begin{remark}[The substrate is sub-political, not anti-political]
The ungoverned substrate is not a rejection of politics or of
governance; it is a layer beneath them. The political question
of how a network state should be governed, the institutional
question of how a DAO should make decisions, the legal
question of how an arrangement should integrate with
conventional jurisdictions --- all of these remain real and
unresolved at the layer at which they apply. The substrate
neither answers nor displaces them. What the substrate does is
make it possible for these questions to be answered
\emph{differently in different arrangements} on the same
infrastructure, rather than uniformly across all arrangements
that share the infrastructure.
\end{remark}

% ════════════════════════════════════════════════════════════════════
\section{Where Governance Returns: At the Graph}
\label{sec:graph-governance}

The argument that the substrate is ungoverned is not the
argument that everything constituted on it is ungoverned. To the
contrary: most useful arrangements constituted on the substrate
will have governance of some kind. The present section
identifies where governance returns and what shape it takes.

\paragraph{Assemblies as governed entities.}
A bonded process tree composed at the substrate is an
assembly --- the corpus's term for a transaction-scoped
arrangement of bonded commitments among multiple parties. A
simple assembly (two parties, one commitment) requires no
governance: the parties have agreed to the commitment, and the
substrate enforces it. A complex assembly (many parties, many
commitments, ongoing relationships) typically does require
governance: questions arise that the original commitments did
not specify, parties join and leave, parameters need to be
updated, disputes require resolution beyond what bonding can
provide. For complex assemblies, governance is real work that
the assembly's participants must do, and the substrate provides
no scaffolding for it.

This is a feature, not a gap. The substrate's silence on
governance allows assemblies to choose the governance form
appropriate to their function. A cooperative assembly can use
cooperative governance; a corporate assembly can use corporate
governance; a network-state assembly can use whatever
constitutional process the network state has settled on; a
peer-network assembly can use no formal governance and rely on
informal coordination. Each of these is a graph-tier choice
that the substrate admits and on which it takes no position.

\paragraph{The COALA scaffolding extends.}
The legal scaffolding work developed within COALA and adjacent
groups remains relevant and useful for assemblies that need it.
A DAO constituted on the substrate that wishes to enter into
contracts with non-DAO counterparties, hold conventional
property, limit liability for participants, or engage with
regulators on recognized terms, will benefit from the legal
frameworks the COALA program has produced. The scaffolding does
not become unnecessary because the substrate is ungoverned; it
becomes one option among many for assemblies that elect to
participate in conventional legal systems. Other assemblies may
elect not to (purely on-chain coordination, transient
relationships, agent-only arrangements), and the substrate
admits both.

\paragraph{How DAOs use the substrate without governing it.}
A common objection to the ungoverned-substrate framing reads as
follows: DAOs need governance, the bonded primitive is meant for
DAOs to use, therefore the bonded primitive needs a governance
interface for DAOs. The objection mistakes a substrate property
for a runtime requirement. DAOs that use the substrate compose
through the same wallet abstraction every other party uses; the
substrate sees a wallet that signs commitments and posts bond,
not a DAO that needs a governance accommodation. The four
composition patterns this section catalogues ---
multisig-as-buyer, governance-executor-as-buyer, streaming
upstream of commit, bookkeeping from chain state --- are the
operational shapes through which a DAO's governance \emph{at the
graph tier} reaches the substrate without requiring the substrate
to take any position on what DAO governance should be.

\emph{Multisig as buyer.} A DAO whose treasury sits in a Safe-style
multisig (or in any threshold-signature wallet) commits as the
buyer using the multisig's standard approval flow. EIP-1271
\citep{eip1271} enables the kernel to verify the multisig's
smart-contract-wallet signature against the same EIP-712 typed
data the protocol uses for EOA signers. The DAO's choice of
signer set, threshold, and rotation policy is a runtime-tier
choice the DAO makes; the substrate is silent on it.

\emph{Governance executor as buyer.} A DAO whose commitments
require formal governance approval --- a quadratic-vote DAO, a
delegation-based DAO, an optimistic-governance DAO --- routes the
authorized commitment through an executor contract that calls
\texttt{commit} after the vote passes
\citep{compound_governor_2020, openzeppelin_governor_2024}. The
substrate sees the executor's signature; the substrate does not
see the vote, the proposal text, the delegation graph, or any
other governance machinery. The DAO's governance is real and
operates entirely above the substrate; the substrate is invisible
to the governance and the governance is invisible to the
substrate.

\emph{Streaming upstream of commit.} A DAO with an ongoing
disbursement program (a grants program, a continuous-funding
public-goods round, a recurring contributor-payment commitment)
can compose a streaming contract upstream of the substrate
\citep{sablier_2023}. The streaming layer handles the
continuous-flow allocation under the DAO's chosen disbursement
policy; the substrate handles the discrete bonded commitments
that disburse the streamed funds against specific work. The two
layers compose without either's accounting model bending.

\emph{Bookkeeping from chain state.} A DAO's accounting and
audit reading work directly from chain state. The DAO's treasury
position, disbursement history, and commitment portfolio for the
disbursed-via-substrate portion of its activity are derivable
from \texttt{OrderCommitted} and \texttt{OrderResolved} events
without any substrate-specific accommodation. This is the same
property the companion accounting paper develops more fully; from
the present paper's perspective the property matters because it
means a DAO using the substrate produces verifiable accounting
records as a byproduct of operation, without the substrate
imposing any accounting interface on the DAO.

The four patterns share a structural property: each compose with
the substrate by treating it as a wallet-and-signature interface
rather than as a governance interface. The substrate provides the
kernel-layer enforcement; the DAO's graph-tier governance is
unaffected and is structurally compatible with whatever
governance shape the DAO has chosen. This is the operational form
of the section's general argument: governance returns at the
graph; the substrate stays out of it; and the standard
institutional shapes already in use among DAO and treasury
operators \citep{messari_dao_treasury_2023, deepdao_treasury_2024}
compose with the substrate without modification.

\paragraph{Network states on the substrate.}
A network-state community constituted on the substrate has its
governance work fully in front of it. The substrate provides
the missing economic engine the network-state program identified
as under-specified, but it provides nothing else: not the
constitutional process, not the community self-conception, not
the diplomatic framework, not the internal adjudication. These
are the network state's to design. What the substrate does is
make those designs possible without requiring the network state
to also design and govern its own commercial-enforcement
infrastructure, which has historically been a substantial
fraction of what state-formation involves.

\paragraph{The infrastructural-literacy concern.}
The blockchain-governance program's concern about
infrastructural literacy --- that participants generally cannot
evaluate the arrangements they are operating under --- applies
to the substrate as it does to any other infrastructure. The
substrate does not reduce the literacy gap; what it does is
\emph{not expand} the gap. A governance-rich protocol's literacy
demand grows over time as the protocol's rules and parameters
evolve; a participant must track the evolution to remain
informed. An invariant kernel's literacy demand is bounded: the
participant must understand the kernel's fixed rules once and
for all. This does not solve the literacy problem (the kernel's
rules are non-trivial, and most participants will not in fact
read them), but it bounds the problem in a way that
governance-rich protocols structurally cannot. The literacy work
the blockchain-governance program calls for is necessary; the
substrate makes it possible without requiring participants to
re-do it on a continuing basis.

\begin{remark}[On the relationship to confidence-machine analysis]
The confidence-machine framing of the
blockchain-governance program correctly identifies that trust
shifted from human institutions to code does not eliminate
trust or governance; it relocates them. The present paper
agrees with this for the class of arrangements the framing
addresses (governance-rich protocols whose rules evolve). The
ungoverned substrate is a different case: trust in the
substrate is required only in the substrate's invariant
specification, not in any ongoing human governance of it. The
trust is therefore mathematizable: formal verification, audit,
and the public visibility of the deployed bytecode discharge
the trust burden in a way that is not available for
governance-rich systems whose evolving parameters cannot be
verified in advance. The confidence-machine concern about
infrastructural literacy survives; the concern about
ongoing-governance accountability does not, because there is no
ongoing governance of the substrate to hold accountable.
\end{remark}

% ════════════════════════════════════════════════════════════════════
\section{What This Paper Is Not}
\label{sec:not}

The argument is liable to several misreadings. We address them
explicitly.

\paragraph{Not a critique of the network-state program.}
The paper does not argue that the network-state program is
wrong, that network states are undesirable, or that the
program's prescriptions should be rejected. The program's
analysis of exit, of community formation, of what it takes to
constitute a sovereign-equivalent entity, is substantive on its
own terms. The paper's claim is narrower: the substrate that
the program needs does not have to be one the network state
constitutes and governs. Recognizing the ungoverned-substrate
option clarifies what the network-state program is committed to
and what is design choice within it; it does not adjudicate
whether the program's overall direction is right.

\paragraph{Not a critique of the blockchain-governance program.}
The paper does not argue that the blockchain-governance program
is wrong, that the confidence-machine analysis is incorrect, or
that COALA's legal-scaffolding work is misguided. The program's
analysis applies to the class of arrangements it was developed
for, and it applies correctly. The paper's claim is that the
class of governance-rich blockchain arrangements is not the only
class; an invariant-kernel architecture is a structurally
distinct class to which the confidence-machine concern about
ongoing governance does not apply, and to which the COALA
scaffolding becomes one option among many for assemblies built
on top of the substrate. This is an extension of the program's
analytical landscape, not a refutation of it.

\paragraph{Not a claim that all blockchain protocols should be ungoverned.}
The paper does not argue that governance-rich protocols are
mistaken, that adaptability through governance is unnecessary,
or that all on-chain coordination should adopt the
ungoverned-substrate model. Many protocols have functions that
genuinely require adaptability and therefore genuinely require
governance: the paper's argument is that the
\emph{coordination-substrate function} is not in this class
when it is restricted to enforcing bonded bilateral commitments,
and that this restriction is what makes the
ungoverned-substrate option available for that specific
function. Protocols whose function is genuinely adaptive should
have governance; the paper does not contest this.

\paragraph{Not a depoliticization of the network question.}
The paper does not argue that politics is over, that governance
has been transcended, or that the substrate change makes
political questions irrelevant. Quite the opposite: the
substrate change is the precondition under which political
questions about network arrangements become legible as
graph-tier questions rather than being entangled with
substrate-tier governance dependencies. There is more
political work to do at the graph layer, not less, because
arrangements that previously had to be uniform across the
infrastructure can now diverge substantively from one another.
The work of designing graph-tier governance for the
arrangements that need it is not addressed by the substrate
change; it is enabled by it.

\paragraph{Not a constitutional proposal.}
The paper does not propose any constitution, governance form,
or political arrangement for any specific network. The
substrate admits arbitrary arrangements; choosing among them is
work that the participants of each arrangement must do. The
paper's contribution is structural: identifying that the
substrate exists as an option, that it is distinct from the
governance-at-substrate-layer assumption that contemporary
programs share, and that the relocation of governance to the
graph is the political-architectural content of recognizing the
option.

\paragraph{Not a prediction.}
Finally, the paper does not predict that the network-state
program will adopt the ungoverned-substrate option, that
blockchain-governance researchers will modify their analyses to
include it, or that any specific arrangement will be built. The
argument is conceptual: the option exists, it is distinct from
what either program currently presupposes, and recognizing it
clarifies the analytical landscape. Whether the recognition
occurs and what arrangements emerge from it are empirical
questions on which the present paper takes no position.

% ════════════════════════════════════════════════════════════════════
\section{Conclusion}
\label{sec:conclusion}

This paper has argued that two contemporary research programs
on network-native coordination --- the network-state program
advanced by \citet{srinivasan2022network} and the
blockchain-governance program developed by
\citet{defilippi_wright_2018} and the COALA-affiliated research
community --- share a structural assumption that
the alternative-to-state coordination must itself be a governed
network entity. The bonded settlement primitive examined in the
corpus's companion mechanism paper realizes a third option that
neither program names: an ungoverned substrate, invariant by
design, on which arbitrary governed entities can be composed but
which itself is not a governed entity at all. Recognizing the
option does not adjudicate among governance forms above the
substrate; it relocates the governance question from the
substrate, where both programs place it, to the graph, where
the substrate admits it to be answered in arbitrarily many
ways.

The substantive observation is that the choice between a
governance-rich substrate and an ungoverned substrate is a real
design choice with real consequences. Governance-rich substrates
provide adaptability at the cost of ongoing
governance-of-the-governance burden; ungoverned substrates
provide invariance at the cost of inability to adapt. For
substrate functions that are amenable to fixed specification ---
the bonded settlement primitive being the case the corpus has
established --- the ungoverned-substrate option is available, and
its availability is what permits the diversity of higher-layer
governance forms the substrate admits. For substrate functions
that are not amenable to fixed specification, the
governance-rich option remains the only one available, and the
blockchain-governance program's analysis of those arrangements
remains directly applicable.

We close with two observations of method, mirroring the closings
of the companion papers in the political-philosophy and
political-economy clusters. First, the kernel is ideologically
agnostic; the graph is the politics --- and the present paper
extends this formulation in a fourth direction. The companion
political-economy papers showed that subordination and
coordination rent become graph-tier choices when the substrate
removes them as load-bearing variables. The companion
political-philosophy paper on coercion showed that coercion
becomes a bounded political-philosophical question when the
substrate substitutes pre-commitment for retributive
enforcement. The present paper shows that governance becomes a
graph-tier choice when the substrate is invariant. In each
case, the move is the same: a feature that prior architectures
treated as constitutive of the substrate becomes a graph-tier
choice when the substrate is designed to exclude it. The graph
is the politics in a strong sense: it is the layer at which
political-architectural choices are made on infrastructure that
makes no political-architectural choices itself.

Second, the recursive note. Future architectures may make
further substrate variables relocatable to the graph, or may
re-introduce variables the present substrate has excluded. The
identification of the ungoverned-substrate option is not a
terminal architectural achievement; it is the next architectural
choice whose conditions and consequences the present moment has
made legible. Subsequent moments will face their own choices,
on conditions the present paper cannot anticipate. What the
substrate change does is enlarge the menu of architectural
options visible from the present, not exhaust the architectural
options that subsequent moments will have available.

\section*{Acknowledgments}
The argument of this paper began as a proposal to extend the
companion political-philosophy paper on coercion to engage two
specific contemporary programs by name. The Network State
program was already in scope as a remark in that paper; the
suggestion to engage Primavera De Filippi and the COALA-aligned
research community substantively was the prompt that occasioned
the present essay. The COALA collaborators were already
acknowledged in the companion political-economy paper on
hegemony; the present engagement is meant in the spirit of that
acknowledgment, as collaboration with rigorous adjacent work
that addresses the same architectural question from a different
direction. I thank the interlocutor whose suggestion prompted
this paper, and the readers of earlier drafts who helped me
calibrate the engagement carefully --- the paper is meant
neither to absorb the adjacent programs into the corpus's
framing nor to dismiss them, but to identify where the corpus's
framing supplies an architectural option the adjacent programs
have not yet had occasion to name.


% ════════════════════════════════════════════════════════════════════
% References
% ════════════════════════════════════════════════════════════════════

\begin{thebibliography}{99}

\bibitem[De~Filippi and Hassan(2016)]{defilippi_hassan_2016}
De~Filippi, P. and Hassan, S.
\newblock Blockchain Technology as a Regulatory Technology: From
  Code is Law to Law is Code.
\newblock \emph{First Monday}, 21(12), 2016.

\bibitem[De~Filippi et al.(2022)]{defilippi_mannan_reijers_2022}
De~Filippi, P., Mannan, M., and Reijers, W.
\newblock Blockchain as a Confidence Machine: The Problem of
  Trust \& Challenges of Governance.
\newblock \emph{Technology in Society}, 62:101284, 2020.
\newblock Extended treatment in subsequent work, 2022.

\bibitem[De~Filippi and Wright(2018)]{defilippi_wright_2018}
De~Filippi, P. and Wright, A.
\newblock \emph{Blockchain and the Law: The Rule of Code}.
\newblock Harvard University Press, Cambridge, MA, 2018.

\bibitem[Hirschman(1970)]{hirschman1970exit}
Hirschman, A.~O.
\newblock \emph{Exit, Voice, and Loyalty: Responses to Decline in
  Firms, Organizations, and States}.
\newblock Harvard University Press, Cambridge, MA, 1970.

\bibitem[Lessig(2006)]{lessig2006code}
Lessig, L.
\newblock \emph{Code: Version 2.0}.
\newblock Basic Books, New York, 2006.

\bibitem[Reijers and Coeckelbergh(2018)]{reijers_coeckelbergh_2018}
Reijers, W. and Coeckelbergh, M.
\newblock The Blockchain as a Narrative Technology: Investigating
  the Social Ontology and Normative Configurations of
  Cryptocurrencies.
\newblock \emph{Philosophy \& Technology}, 31(1):103--130, 2018.

\bibitem[Walch(2017)]{walch2017path}
Walch, A.
\newblock The Path of the Blockchain Lexicon (and the Law).
\newblock \emph{Boston University Review of Banking and Financial
  Law}, 36:713--765, 2017.

\bibitem[Walch(2019)]{walch2019deconstructing}
Walch, A.
\newblock Deconstructing `Decentralization': Exploring the Core
  Claim of Crypto Systems.
\newblock In C.~Brummer, editor, \emph{Crypto-Assets: Legal and
  Regulatory Implications of Distributed Ledger Technology}.
  Oxford University Press, 2019.

\bibitem[Werbach(2018)]{werbach2018blockchain}
Werbach, K.
\newblock \emph{The Blockchain and the New Architecture of Trust}.
\newblock MIT Press, Cambridge, MA, 2018.

\bibitem[ERC-1271(2018)]{eip1271}
Giordano, F., Lockyer, M., and Castonguay, P.
\newblock EIP-1271: Standard Signature Validation Method for
  Contracts.
\newblock Ethereum Improvement Proposal 1271, July 2018.
\newblock \url{https://eips.ethereum.org/EIPS/eip-1271}

\bibitem[Compound(2020)]{compound_governor_2020}
Compound Labs.
\newblock \emph{Compound Governor: On-Chain Governance Module}.
\newblock Compound Protocol Documentation, 2020.

\bibitem[OpenZeppelin(2024)]{openzeppelin_governor_2024}
OpenZeppelin.
\newblock \emph{Governor Contracts: Reference Implementation for
  On-Chain Governance}.
\newblock OpenZeppelin Contracts Documentation, ongoing as of 2024.

\bibitem[Sablier(2023)]{sablier_2023}
Sablier Labs.
\newblock \emph{Sablier Protocol: Token Streaming Infrastructure}.
\newblock Sablier Documentation, 2023.

\bibitem[Messari(2023)]{messari_dao_treasury_2023}
Messari Research.
\newblock \emph{State of DAO Treasuries 2023}.
\newblock Messari Quarterly Report, 2023.

\bibitem[DeepDAO(2024)]{deepdao_treasury_2024}
DeepDAO.
\newblock \emph{DAO Treasury Aggregate Statistics}.
\newblock DeepDAO Analytics Platform, ongoing as of 2024.

\bibitem[Srinivasan(2022)]{srinivasan2022network}
Srinivasan, B.
\newblock \emph{The Network State: How to Start a New Country}.
\newblock Self-published, 2022.

\end{thebibliography}

\end{document}
